% GENERATED FILE. Edit canonical lessons, then run npm run generate:books.
% canonical-source-hash: fnv1a64:a0a98bc067f53be9
% canonical-lessons: MR-C464-himmat, MR-C464-jidda, MR-C464-cikati, MR-C464-umed, MR-C464-tag-dharne

\chapter{Courage and Grit}
\label{ch:mr-courage-and-grit}
\hlchaptermodality{464}

\begin{chapteropening}
\textbf{By the end of this chapter:} \emph{I can talk about nerve, grit, zest and holding out.}

Complete the last lesson of chapter 464: I can talk about nerve, grit, zest and holding out.
\end{chapteropening}

\section[\texorpdfstring{\mr{हिंमत} (himmat)}{हिंमत (himmat)}]{\mr{हिंमत} (himmat) — nerve; courage}

\label{lesson:MR-C464-himmat}



\begin{quote}
Before the new one: say the Marathi for to feel ashamed, then the Marathi for to feel sheepish.
\end{quote}

\subsection*{You'll want to know: \mr{हिंमत}}
\textbf{\mr{हिंमत}} — \emph{himmat} — \textquotedblleft{}nerve; courage\textquotedblright{}.

\begin{quote}
\mr{रात्री} \mr{एकट्याने} \mr{जायची} \mr{त्याची} \mr{हिंमत} \mr{झाली} \mr{नाही}.
\end{quote}

\begin{quote}
\emph{rātrī ekṭyāne jāycī tyācī himmat jhālī nāhī.}
\end{quote}

\begin{quote}
He did not have the nerve to go alone at night.
\end{quote}

Feminine.

\mr{हिंमत} \mr{होणे}, to dare, with the genitive: \mr{माझी} \mr{हिंमत} \mr{झाली}. \mr{हिंमत} \mr{हरणे}, to lose heart.

\begin{cousinweb}
\textbf{\mr{हिंमत}} (\emph{himmat}) is a loan from Persian \emph{himmat}, usually traced to Arabic \emph{himma}, \textquotedblleft{}resolve\textquotedblright{}.
\end{cousinweb}

\begin{grammarlens}[title={What you've built}]
One more word for describing what you lived and felt.
\end{grammarlens}

\subsection*{Your turn}
\begin{itemize}
\raggedright
  \item \emph{Say it:} \emph{himmat}
  \item \emph{Say it:} \emph{himmat}, once more
  \item \emph{Say it:} \emph{rātrī ekṭyāne jāycī tyācī himmat jhālī nāhī.} — He did not have the nerve to go alone at night.
  \item \emph{Recall:} say the Marathi for to feel ashamed, then the Marathi for to feel sheepish, then say \emph{himmat} again
\end{itemize}

\begin{tcolorbox}[breakable,colback=teal!4,colframe=teal!35!black,title={Before you move on}]
What does \mr{हिंमत} mean? (\textquotedblleft{}Nerve; courage\textquotedblright{}.) Say it once more.
\end{tcolorbox}
\section[\texorpdfstring{\mr{जिद्द} (jidda)}{जिद्द (jidda)}]{\mr{जिद्द} (jidda) — grit; tenacity}

\label{lesson:MR-C464-jidda}



\begin{quote}
Before the new one: say the Marathi for to feel sheepish, then the Marathi for nerve.
\end{quote}

\subsection*{You'll want to know: \mr{जिद्द}}
\textbf{\mr{जिद्द}} — \emph{jidda} — \textquotedblleft{}grit; tenacity\textquotedblright{}.

\begin{quote}
\mr{जिद्दीने} \mr{अभ्यास} \mr{करून} \mr{ती} \mr{उत्तीर्ण} \mr{झाली}.
\end{quote}

\begin{quote}
\emph{jiddīne abhyās karūn tī uttīrṇa jhālī.}
\end{quote}

\begin{quote}
She studied with grit and passed.
\end{quote}

Feminine.

\mr{जिद्दीने}, doggedly; \mr{जिद्दी}, tenacious.

\begin{cousinweb}
\textbf{\mr{जिद्द}} (\emph{jidda}) is a loan from Hindustani \emph{zidd}, \textquotedblleft{}insistence\textquotedblright{}, usually traced to Arabic \emph{didd}, \textquotedblleft{}opposite\textquotedblright{}; Marathi writes the \emph{z} as \mr{ज}.
\end{cousinweb}

\begin{grammarlens}[title={What you've built}]
One more word for describing what you lived and felt.
\end{grammarlens}

\subsection*{Your turn}
\begin{itemize}
\raggedright
  \item \emph{Say it:} \emph{jidda}
  \item \emph{Say it:} \emph{jidda}, once more
  \item \emph{Say it:} \emph{jiddīne abhyās karūn tī uttīrṇa jhālī.} — She studied with grit and passed.
  \item \emph{Recall:} say the Marathi for to feel sheepish, then the Marathi for nerve, then say \emph{jidda} again
\end{itemize}

\begin{tcolorbox}[breakable,colback=teal!4,colframe=teal!35!black,title={Before you move on}]
What does \mr{जिद्द} mean? (\textquotedblleft{}Grit; tenacity\textquotedblright{}.) Say it once more.
\end{tcolorbox}
\section[\texorpdfstring{\mr{चिकाटी} (cikāṭī)}{चिकाटी (cikāṭī)}]{\mr{चिकाटी} (cikāṭī) — perseverance}

\label{lesson:MR-C464-cikati}



\begin{quote}
Before the new one: say the Marathi for nerve, then the Marathi for grit.
\end{quote}

\subsection*{You'll want to know: \mr{चिकाटी}}
\textbf{\mr{चिकाटी}} — \emph{cikāṭī} — \textquotedblleft{}perseverance\textquotedblright{}.

\begin{quote}
\mr{चिकाटीने} \mr{सराव} \mr{केला} \mr{की} \mr{यश} \mr{मिळतं}.
\end{quote}

\begin{quote}
\emph{cikāṭīne sarāv kelā kī yaś miḷtaṁ.}
\end{quote}

\begin{quote}
If you practise with perseverance, success comes.
\end{quote}

Feminine.

Feminine; \mr{चिकाटी} \mr{सोडणे}, to give up.

\begin{cousinweb}
\textbf{\mr{चिकाटी}} (\emph{cikāṭī}) is a Marathi noun whose origin is not settled; it is often linked with \mr{चिकट}, \textquotedblleft{}sticky\textquotedblright{}.
\end{cousinweb}

\begin{grammarlens}[title={What you've built}]
One more word for describing what you lived and felt.
\end{grammarlens}

\subsection*{Your turn}
\begin{itemize}
\raggedright
  \item \emph{Say it:} \emph{cikāṭī}
  \item \emph{Say it:} \emph{cikāṭī}, once more
  \item \emph{Say it:} \emph{cikāṭīne sarāv kelā kī yaś miḷtaṁ.} — If you practise with perseverance, success comes.
  \item \emph{Recall:} say the Marathi for nerve, then the Marathi for grit, then say \emph{cikāṭī} again
\end{itemize}

\begin{tcolorbox}[breakable,colback=teal!4,colframe=teal!35!black,title={Before you move on}]
What does \mr{चिकाटी} mean? (\textquotedblleft{}Perseverance\textquotedblright{}.) Say it once more.
\end{tcolorbox}
\section[\texorpdfstring{\mr{उमेद} (umed)}{उमेद (umed)}]{\mr{उमेद} (umed) — zest; hope}

\label{lesson:MR-C464-umed}



\begin{quote}
Before the new one: say the Marathi for grit, then the Marathi for perseverance.
\end{quote}

\subsection*{You'll want to know: \mr{उमेद}}
\textbf{\mr{उमेद}} — \emph{umed} — \textquotedblleft{}zest; hope\textquotedblright{}.

\begin{quote}
\mr{आजोबा} \mr{म्हातारे} \mr{झाले} \mr{तरी} \mr{त्यांची} \mr{उमेद} \mr{कायम} \mr{आहे}.
\end{quote}

\begin{quote}
\emph{ājobā mhātāre jhāle tarī tyāñcī umed kāyam āhe.}
\end{quote}

\begin{quote}
Grandfather has grown old, but his zest is still there.
\end{quote}

Feminine.

\mr{उमेद} \mr{हरणे}, to lose heart; \mr{उमेदीचे} \mr{दिवस}, the years of youthful drive.

\begin{cousinweb}
\textbf{\mr{उमेद}} (\emph{umed}) is a loan from Persian \emph{umed}, \textquotedblleft{}hope\textquotedblright{}.
\end{cousinweb}

\begin{grammarlens}[title={What you've built}]
One more word for describing what you lived and felt.
\end{grammarlens}

\subsection*{Your turn}
\begin{itemize}
\raggedright
  \item \emph{Say it:} \emph{umed}
  \item \emph{Say it:} \emph{umed}, once more
  \item \emph{Say it:} \emph{ājobā mhātāre jhāle tarī tyāñcī umed kāyam āhe.} — Grandfather has grown old, but his zest is still there.
  \item \emph{Recall:} say the Marathi for grit, then the Marathi for perseverance, then say \emph{umed} again
\end{itemize}

\begin{tcolorbox}[breakable,colback=teal!4,colframe=teal!35!black,title={Before you move on}]
What does \mr{उमेद} mean? (\textquotedblleft{}Zest; hope\textquotedblright{}.) Say it once more.
\end{tcolorbox}
\section[\texorpdfstring{\mr{तग} \mr{धरणे} (tag dharṇe)}{तग धरणे (tag dharṇe)}]{\mr{तग} \mr{धरणे} (tag dharṇe) — to hold out}

\label{lesson:MR-C464-tag-dharne}



\begin{quote}
Before the new one: say the Marathi for perseverance, then the Marathi for zest.
\end{quote}

\subsection*{You'll want to know: \mr{तग} \mr{धरणे}}
\textbf{\mr{तग} \mr{धरणे}} — \emph{tag dharṇe} — \textquotedblleft{}to hold out\textquotedblright{}.

\begin{quote}
\mr{दुष्काळातही} \mr{हे} \mr{झाड} \mr{तग} \mr{धरून} \mr{उभं} \mr{आहे}.
\end{quote}

\begin{quote}
\emph{duṣkāḷātahī he jhāḍ tag dharūn ubhaṁ āhe.}
\end{quote}

\begin{quote}
Even in the drought this tree has held out.
\end{quote}

Past, with -\mr{ने} (he …): \mr{त्याने} \mr{तग} \mr{धरला}, \emph{tyāne tag dharlā}.

To hold one's ground through hardship: \mr{तग} \mr{धरून} \mr{राहणे}.

\begin{cousinweb}
\textbf{\mr{तग} \mr{धरणे}} (\emph{tag dharṇe}) is a Marathi phrase built on \mr{तग}, from the reconstructed \emph{*tagg}, \textquotedblleft{}endure\textquotedblright{}, under which CDIAL lists Marathi \emph{tagṇẽ}, \textquotedblleft{}to last\textquotedblright{} (CDIAL 5625), with \mr{धरणे}, \textquotedblleft{}to hold\textquotedblright{}.
\end{cousinweb}

\begin{grammarlens}[title={What you've built}]
One more word for describing what you lived and felt.
\end{grammarlens}

\subsection*{Your turn}
\begin{itemize}
\raggedright
  \item \emph{Say it:} \emph{tag dharṇe}
  \item \emph{Say it:} \emph{tag dharṇe}, once more
  \item \emph{Say it:} \emph{duṣkāḷātahī he jhāḍ tag dharūn ubhaṁ āhe.} — Even in the drought this tree has held out.
  \item \emph{Recall:} say the Marathi for perseverance, then the Marathi for zest, then say \emph{tag dharṇe} again
\end{itemize}

\begin{tcolorbox}[breakable,colback=teal!4,colframe=teal!35!black,title={Before you move on}]
What does \mr{तग} \mr{धरणे} mean? (\textquotedblleft{}To hold out\textquotedblright{}.) Say it once more.
\end{tcolorbox}
